\subsection{Sustitución del prefactor conservando los términos de orden
\texorpdfstring{$1/k$}{1/k}}\label{sec:B4T07}
Se usan $S_0,\Gamma$ de la Sección~\ref{sec:B4T06} y se define
\[
 m_{2,a}=\sum_{i\ni a}p_i^2,\quad
 \kappa_i=a_i^T\Gamma a_i-4=O(k^{-1}).
\]
En un soporte fijo de la clase cuantitativa,
\begin{equation}\label{eq:ell-exp}
 \boxed{\begin{aligned}
 \ell={}&-\tfrac12\log\det S_{0,\rm gf}-3\sum p_i^2-14\sum p_i^3
 -\tfrac32\sum p_i^2\kappa_i
 +\tfrac94\sum_a m_{2,a}^2+O(k^{-2}).
 \end{aligned}}
\end{equation}
\paragraph{Demostración de la expansión.}
Se tiene $S-S_0=A^T\operatorname{diag}(3p^2+12p^3+O(p^4))A=E$,
con $\|E\|_\infty=O(k^{-1})$, diagonal $3m_{2,a}+O(k^{-2})$
y entradas fuera de la diagonal $O(k^{-2})$.
En el complemento del núcleo, el resto cúbico del logaritmo del determinante es $O(k^{-2})$.
El término lineal es
\[
 \operatorname{tr}(\Gamma E)=12\sum p^2+
 3\sum p^2\kappa+48\sum p^3+O(k^{-2}).
\]
El término cuadrático es
$\operatorname{tr}(\Gamma E\Gamma E)=9\sum_a m_{2,a}^2+O(k^{-2})$.
Para comprobarlo se usan $\Gamma=I+O_{\max}(k^{-1})$, con norma $\infty$
acotada, y la estructura de $E$ en la diagonal y fuera de ella.
Además $-\tfrac12\log b(p)=3p^2+10p^3+O(p^4)$ y
$\sum p^4\le(C/k)^3\sum p=O(k^{-2})$.
Esto demuestra todos los coeficientes, incluido $-14$.

\paragraph{Cambios de la pseudoinversa en líneas supervivientes.}
Sean $\Gamma'$ la pseudoinversa de $S_{0,\rm hijo}$ en $p^*_{\rm hijo}$ y
$\Delta_0=S_{0,\rm hijo}-S_{0,II}$.
La diagonal de $\Delta_0$ es cero, y su norma $\infty$ es $O(k^{-1})$.
El complemento de Schur parental tiene inverso
$\Pi_I^\perp\Gamma_{II}\Pi_I^\perp=\Gamma_{II}+O_{\max}(k^{-2})$:
se usa $\Gamma_{II}v_I=-\Gamma_{IJ}v_J=O_\infty(k^{-1})$.
El término de Schur de rango cuatro tiene entradas $O(k^{-2})$.
Aplicando la identidad de la resolvente en el complemento de $v_I$, se obtiene
\[
 \boxed{\Gamma'_{ab}-\Gamma_{ab}=-(\Delta_0)_{ab}+O(k^{-2}),\quad a,b\in I.}
\]
Para ver el resto por entradas, obsérvese que la norma $\infty$ de cada factor
perturbador es $O(k^{-1})$; los productos cuadráticos son $O(k^{-2})$.
En el término lineal, sustituir cada inversa por $I$ produce un cambio $O(k^{-2})$ por entrada.
Si $a,b$ son dos líneas de una casilla superviviente $i$, su única
intersección es $i$, de modo que $(\Delta_0)_{ab}=y_i=O(k^{-2})$.
Así $\kappa'_i-\kappa_i=O(k^{-2})$ en $E'$.

Los términos $\sum p^3$, $\sum p^2\kappa$ y $\sum_a m_{2,a}^2$
cambian cada uno en $O(k^{-2})$: hay $O(k)$ casillas eliminadas;
en las supervivientes, $y/p=O(k^{-1})$; por línea retenida,
$m'_{2,a}-m_{2,a}=O(k^{-2})$; y se eliminan sólo cuatro líneas.
Por tanto
\begin{equation}\label{eq:ell-hijo}
 \boxed{\ell_{\rm hijo}-\ell_{\rm padre}
 =-\tfrac12\log\frac{\det S_{0,\rm hijo,gf}}{\det S_{0,\rm padre,gf}}
 -3\left[\sum_{E'}(p_i+y_i)^2-\sum_Ep_i^2\right]+O(k^{-2}).}
\end{equation}
La sustitución no consiste en omitir un error global $O(k^{-1})$:
se han identificado sus términos y se ha probado que su cambio es $O(k^{-2})$.

Reunimos lo probado en esta sección.
\begin{lemasub}[Sustitución del prefactor]\label{lem:prefactor}
En un soporte fijo de la clase cuantitativa vale~\eqref{eq:ell-exp}. En las hipótesis de la
Sección~\ref{sec:B4T01}, con $y=p^*_{\rm hijo}-w$, vale además~\eqref{eq:ell-hijo}.
\end{lemasub}
